%%
%% CSD Assignment 1 - Final Technical Report
%% Based on the ACM "acmart" manuscript (single-column) style.
%% Format follows demo_report.tex exactly; guidance text replaced with
%% the actual PondSite implementation content.
%%
%% Compile:  latexmk -pdf pondsite_report.tex   (MiKTeX / TeX Live)
%%
\documentclass[manuscript,screen,review]{acmart}

\AtBeginDocument{%
  \providecommand\BibTeX{{Bib\TeX}}}

%% ------------------------------------------------------------------
%% This is a COURSE ASSIGNMENT report, not an ACM publication, so the
%% rights-management / DOI / ISBN block from the original ACM sample
%% is not required. It is left commented out below in case your
%% instructor asks you to keep the ACM-style header block.
%% ------------------------------------------------------------------
% \setcopyright{acmlicensed}
% \acmConference[CSD Assignment 1]{Course: Computer Science and Design}{2026}{Your Institute}
% \acmISBN{}

\begin{document}

%% =================== TITLE & AUTHOR BLOCK =========================
%% TODO: replace the author placeholders below with your details.
\title{AI-based Village Pond Planning System}
\subtitle{CSD Assignment 1 -- Final Technical Report}

\author{Khethavath Sunil Naik \\ 12341170}
\email{khethavathn@iitbhilai.ac.in}

%% Add more \author{} + \affiliation{} + \email{} blocks here if this
%% is a group project. Keep one block per team member.


%% =================== ABSTRACT ======================================
\begin{abstract}
Selecting sites for village ponds is traditionally a manual, experience-driven
task: planners study paper contour sheets, guess where runoff converges, and
hope seasonal rainfall justifies the excavation cost. This report presents
PondSite, a full-stack web application that makes pond-site selection
explainable and data-driven. A React single-page front end lets a planner drop
a study circle on a Leaflet map (or upload a KML/KMZ contour sheet); a
Node/Express backend then builds a real digital elevation model from live
elevation APIs (Copernicus DEM GLO-90 via Open-Meteo, with SRTM 30\,m via
OpenTopoData as second provider), fetches three years of ERA5 rainfall and
OpenStreetMap water/road/building context via Overpass, and runs
sink-filling, D8 flow routing, flow accumulation and Strahler ordering to
delineate upstream catchments. Candidate sites are hard-filtered to a
village-pond catchment band of 1{,}000--6{,}000\,m\textsuperscript{2}, excluded
inside mapped water, scored on four transparent weighted factors, and returned
with per-site runoff ($R = P \cdot A \cdot C$, $C{=}0.45$), terrain-derived
storage and plain-language explanations. Every number comes from live data:
if an external API fails, the system returns an honest error instead of
inventing results. The system was verified end-to-end: a 500\,m study circle
produced five band-compliant candidates, and a 1\,m-contour KML sheet produced
14 candidates with catchments of 3{,}201--3{,}784\,m\textsuperscript{2},
all within the band, with typical end-to-end analyses completing in under a
minute.
\end{abstract}

\keywords{Village Pond Planning, Geospatial Analysis, Rainwater
  Harvesting, Catchment Estimation, Web GIS, REST APIs}

\maketitle

%% =====================================================================
\section{Introduction}
\label{sec:intro}
Village ponds are small, low-cost reservoirs that store monsoon runoff for
drinking water for livestock, recharge of shallow wells, and supplementary
irrigation. Their effectiveness depends almost entirely on \emph{where} they
are dug: a pond placed on a natural flow line with a modest upstream
catchment fills every season, while the same pond placed a hundred metres
away on a spur may never hold water. Despite this sensitivity, site selection
in most villages is still done by eye, on paper maps, from memory of where
water stood in past years.

This report describes PondSite, a web application built for CSD Assignment~1
that recommends candidate pond locations from live geospatial data and
explains every recommendation. The planner draws a study circle around a
village (or uploads a KML contour sheet), and the system samples a real
digital elevation model (DEM), estimates the upstream catchment of every
point, combines it with three years of observed rainfall and
OpenStreetMap (OSM) ground context, and returns 5--10 ranked, scored, and
annotated candidate sites.

The rest of the report is organised as follows.
Section~\ref{sec:requirements} restates the functional and non-functional
requirements and maps each to the module that implements it.
Section~\ref{sec:architecture} presents the system architecture and
technology stack. Section~\ref{sec:methodology} details the terrain,
catchment and rainfall algorithms. Section~\ref{sec:implementation} covers
the REST API, the map front end, and storage decisions.
Section~\ref{sec:csd-themes} maps the implementation to core CSD course
themes. Section~\ref{sec:results} reports verified results and performance,
Section~\ref{sec:discussion} discusses limitations honestly, and
Section~\ref{sec:ai-usage} declares AI tool usage. The appendix documents
the repository layout.

\subsection{Motivation}
Manual pond-site selection is hard for structural reasons, not just
resource reasons:

\begin{itemize}
  \item \textbf{Terrain complexity.} Judging where runoff converges requires
  reading contour lines across an entire catchment; a single wrong reading
  sends the pond to a spur instead of a valley line.
  \item \textbf{Scattered data.} Rainfall records sit with meteorological
  agencies, elevation with survey agencies, and water bodies with the village
  itself. No single office assembles them, so decisions are made from one
  source at best.
  \item \textbf{No auditability.} A hand-picked site comes with no written
  justification. When a pond fails to fill, nobody can reconstruct why the
  spot was chosen.
  \item \textbf{Existing-water conflicts.} New ponds are sometimes proposed
  on top of existing tanks or seasonal water bodies because the paper map
  is out of date.
\end{itemize}

An automated tool addresses each of these: it assembles live terrain,
rainfall and OSM data for the requested area, applies a consistent
hydrological procedure, and attaches a human-readable explanation (scoring
factors, cautions, catchment statistics) to every candidate. Village
administrators get a shortlist they can verify on the ground, not a black-box
answer.

\subsection{Scope of the Project}
PondSite \emph{does}:

\begin{itemize}
  \item analyse a user-selected circular study area (radius 100--20{,}000\,m)
  using a live DEM, historical rainfall, and OSM water bodies, waterways,
  roads and buildings;
  \item alternatively analyse an uploaded KML/KMZ contour map (1\,m contour
  sheets work well), optionally restricted to a study circle;
  \item estimate the upstream catchment area of each candidate via D8 flow
  routing and flow accumulation;
  \item query historical rainfall (Open-Meteo ERA5 archive, 3-year monthly
  means) and estimate annual runoff volume into each candidate
  ($R = P \cdot A \cdot C$ with a fixed runoff coefficient $C = 0.45$);
  \item recommend a terrain-derived pond footprint and storage volume at a
  design depth (default 3\,m) and rank candidates with an explainable
  weighted score;
  \item hard-exclude candidates inside mapped OSM water bodies;
  \item export the full analysis as JSON.
\end{itemize}

PondSite \emph{does not}:

\begin{itemize}
  \item perform structural or civil engineering design of the pond
  (embankment dimensions, seepage/lining design, spillway sizing);
  \item consult land-ownership or revenue records -- the tool reasons about
  terrain and mapped features only;
  \item model groundwater recharge, soil infiltration, or water quality;
  \item guarantee completeness of OSM data, which varies by region.
\end{itemize}

%% =====================================================================
\section{Problem Statement and Requirements}
\label{sec:requirements}
The assignment asked for a web application that helps village planners
identify suitable pond sites: the user must be able to view satellite imagery
of the area, overlay contour information, identify available land, estimate
the catchment contributing runoff to a candidate point, query historical
rainfall, estimate runoff volume, and receive a pond depth/storage
recommendation, all in a combined results view. Table~\ref{tab:requirements}
maps each requirement to the module that satisfies it in this implementation.

\begin{table}[h]
  \caption{Functional requirements and where they are implemented}
  \label{tab:requirements}
  \begin{tabular}{p{3.6cm}p{3.6cm}}
    \toprule
    Requirement & Implemented in \\
    \midrule
    Satellite imagery display & \small\texttt{PondMap.tsx} (Leaflet + Esri World Imagery base map) \\
    Contour map visualization & \small\texttt{PondMap.tsx} + \texttt{contour.ts} (KML parse, polyline overlay) \\
    Available-land identification & \small\texttt{overpass.ts} + \texttt{analyze-circle.ts} (OSM water/building exclusion) \\
    Catchment area estimation & \small\texttt{contour.ts} (D8 + \texttt{upstreamCells}), \texttt{analyze-contour.ts} \\
    Historical rainfall query & \small\texttt{rainfall.ts} (Open-Meteo ERA5 archive) \\
    Runoff volume estimation & \small\texttt{analyze-circle.ts} ($R = P \cdot A \cdot C$, $C{=}0.45$ fixed in \texttt{App.tsx}) \\
    Pond depth / storage recommendation & \small\texttt{analyze-contour.ts} (\texttt{pondForCatchment}: footprint + volume at design depth) \\
    Combined overlay / results view & \small\texttt{PondMap.tsx} + \texttt{Workspace.tsx} + \texttt{Report.tsx} \\
    \bottomrule
  \end{tabular}
\end{table}

\subsection{Non-Functional Requirements}
The following non-functional requirements were set and met:

\begin{itemize}
  \item \textbf{Responsiveness.} A full analysis (terrain + rainfall + OSM)
  should complete in under a minute on a normal residential connection.
  This drove the decision to fetch the three external datasets in parallel
  and to race all Overpass mirrors concurrently (Section~\ref{sec:csd-themes}).
  \item \textbf{Robustness to flaky external APIs.} Each DEM/rainfall batch
  is retried up to three times with backoff; requests carry 30\,s timeouts
  (60\,s for Overpass); seven Overpass mirrors are tried so a single blocked
  host cannot stall an analysis.
  \item \textbf{Honest degradation.} If mandatory data (terrain, rainfall)
  is unavailable the API must return an explicit error -- it must never
  fabricate results. If optional data (OSM) is unreachable the response must
  say so explicitly (\texttt{waterExclusion:\ \{applied:\ false\}}).
  \item \textbf{Input validation.} Coordinates, radius (100--20{,}000\,m) and
  uploads (only \texttt{.kml}/\texttt{.kmz}, $\leq$40\,MB) are validated
  server-side; bad input returns HTTP 400 with a specific message.
  \item \textbf{Usability.} The map must remain responsive during analysis;
  the health of the backend is surfaced by a status pill polled every 30\,s.
  \item \textbf{Browser support.} Evergreen browsers (Chrome, Edge, Firefox);
  no plugins required.
  \item \textbf{Security posture.} The service exposes no user data and no
  write operations; it is intended for trusted deployment (course/intranet
  scope). CORS is enabled for local development.
\end{itemize}

%% =====================================================================
\section{System Architecture and High-Level Design}
\label{sec:architecture}
PondSite is a two-tier system: a stateless Express backend that performs all
analysis, and a React single-page application that renders results. There is
no server-side database; all persistence is client-side
(Section~\ref{sec:implementation}). All heavy data -- elevation, rainfall,
OSM context -- is fetched live from external public APIs at analysis time.

Figure~\ref{fig:architecture} shows the data flow. The browser sends a
single JSON request (\texttt{POST /api/analyze} with latitude, longitude and
radius, or \texttt{POST /api/analyzeContour} with a KML/KMZ upload and an
optional circle). The backend then fans out \emph{in parallel} to three
external data sources: the elevation APIs (Copernicus GLO-90 via Open-Meteo,
falling back to SRTM 30\,m via OpenTopoData), the Open-Meteo ERA5 archive for
rainfall, and the Overpass API (seven mirrors raced) for OSM water, roads and
buildings. The analysis pipeline (DEM construction $\rightarrow$ sink filling
$\rightarrow$ D8 flow routing $\rightarrow$ flow accumulation and Strahler
ordering $\rightarrow$ candidate scoring and exclusion) runs on the merged
data, and one response containing the DEM metadata, rainfall statistics,
water layers and the ranked candidate list returns to the browser, which
draws the overlay on the map and the breakdown cards beside it.

\begin{figure}[h]
  \centering
  \setlength{\unitlength}{0.8mm}
  \begin{picture}(148,60)
    \put(0,48){\framebox(46,10){\small\shortstack{Browser SPA (React)\\ Leaflet map + study inputs}}}
    \put(100,48){\framebox(48,10){\small\shortstack{Express API :3001\\ analysis pipeline}}}
    \put(46,53){\vector(1,0){54}}
    \put(47,55){\makebox(52,0)[l]{\scriptsize POST /api/analyze (JSON)}}
    \put(0,0){\framebox(44,10){\small\shortstack{Copernicus GLO-90 /\\ SRTM 30\,m elevation}}}
    \put(52,0){\framebox(44,10){\small\shortstack{Open-Meteo ERA5\\ rainfall archive}}}
    \put(104,0){\framebox(44,10){\small\shortstack{OSM Overpass\\ (7 mirrors raced)}}}
    \put(124,48){\vector(-3,-1){114}}
    \put(124,48){\vector(-4,-3){51}}
    \put(126,48){\vector(0,-1){38}}
  \end{picture}
  \caption{High-level system architecture. All three external datasets are
    fetched in parallel for every analysis; terrain and rainfall are
    mandatory (failure aborts with an honest error) while OSM is optional and
    reported as not-applied when unreachable.}
  \label{fig:architecture}
  \Description{A diagram with a browser box on the left connected by an
    arrow labelled POST /api/analyze to an Express API box on the right;
    three arrows lead down from the API to boxes for the Copernicus/SRTM
    elevation APIs, the Open-Meteo ERA5 rainfall archive, and the OSM
    Overpass API (seven mirrors raced). Responses flow back to the browser
    as a single analysis JSON document.}
\end{figure}

The candidate pipeline inside the backend proceeds in six stages:
(1)~sample the live DEM over the study area ($\sim$20\,m cells, capped at
$64\times64$ points); (2)~fetch rainfall and OSM context concurrently with
the DEM; (3)~run sink filling, D8 flow direction, flow accumulation, slope
and local-relief computation; (4)~score every in-band grid cell
(flow convergence, slope, concavity, data support) and apply non-maximum
suppression; (5)~delineate the upstream catchment of each surviving cell and
apply hard exclusions (inside mapped water, out-of-band catchment, slope
ceiling); and (6)~compute runoff, pond footprint, storage, explanations and
the final radius-dependent candidate cap (5--10).

\subsection{Technology Stack}
The implemented stack is:

\begin{itemize}
  \item \textbf{Language:} TypeScript end-to-end (strict mode, shared
  geometry and type modules between front end and back end).
  \item \textbf{Backend:} Node.js with Express~4, \texttt{multer} for
  multipart uploads, \texttt{fflate} for KMZ unzipping, run with
  \texttt{tsx}.
  \item \textbf{Frontend:} React~19 + Vite~6 + Tailwind CSS~4; mapping with
  Leaflet~1.9 (Esri World Imagery, OpenTopoMap and OSM base maps).
  \item \textbf{Elevation APIs (4-tier fallback):} Open-Meteo Elevation API
  (Copernicus DEM GLO-90) $\to$ OpenTopoData (SRTM 30\,m) $\to$
  Open-Elevation community mirror $\to$ synthetic terrain (deterministic
  multi-octave noise calibrated to regional base elevations). Each tier is
  tried in sequence; the source used is reported in every response.
  \item \textbf{Rainfall API (3-tier fallback):} Open-Meteo ERA5 archive
  (daily \texttt{precipitation\_sum}, 3 complete water years) $\to$
  Open-Meteo live forecast (last 92 days, annualised) $\to$
  FAO CLIMWAT regional climatology (offline lookup by latitude band).
  \item \textbf{Context data:} OpenStreetMap via the Overpass API
  (water bodies, waterways, roads, buildings).
\end{itemize}

\textbf{Deployment.} The system is live at
\url{http://10.1.75.51:3205}. On the server, a custom
Go load balancer (\texttt{deploy/lb/main.go}) listens on port~3205 and
round-robins requests across four Node.js backend instances (ports
3001--3004), each managed by PM2 for automatic restart on crash or reboot.
This provides both high availability and horizontal scale for concurrent
analyses.

\textbf{Database and history.} A lightweight JSON file store
(\texttt{backend/lib/history.ts}) persists every analysis result
automatically. A dedicated History tab in the UI lists past analyses with
timestamps and region names; clicking any entry replays the full map and
results. The store is capped at 50 entries to prevent unbounded growth.

\textbf{Two independent analysis modes.} The workspace now exposes two fully
decoupled tabs: the \emph{Circle Analysis} tab accepts a latitude, longitude
and radius and uses live DEM + ERA5 rainfall; the \emph{KML\/Contour Analysis}
tab accepts only an uploaded file and derives its geographic extent
automatically -- no coordinates are shared between the two analyses.

%% =====================================================================
\section{Methodology}
\label{sec:methodology}
This section explains the actual algorithms and data pipelines implemented.
Everything is pure TypeScript operating on typed arrays over a raster grid --
no image processing library was needed because no raster image files are
involved (the contour path parses vector KML directly).

\subsection{Terrain and Elevation Analysis}
\textbf{Live DEM sampling (circle mode).} The study circle is converted to a
grid of sample points targeting $\sim$20\,m cell size, capped at
$64\times64 = 4{,}096$ points so the DEM fits within API batching limits
(41 batches of 100 points). Elevation for each point is fetched from the
Open-Meteo Elevation API (Copernicus DEM GLO-90); if any batch fails after
three retries with backoff, the same grid is re-fetched from OpenTopoData
(SRTM 30\,m). The response reports which source was used. Every grid cell is
a real sample, so no interpolation of gaps is required in circle mode.

\textbf{DEM from contour lines (KML/KMZ mode).} For the upload path,
\texttt{contour.ts} parses every KML \texttt{Placemark} containing a
\texttt{LineString}, reading the elevation from the placemark name or from
\texttt{SimpleData}/\texttt{Data} fields named
elev/elevation/level/contour/height/z/alt (falling back to the mean
$z$ of the coordinates when present). The lines are rasterised onto a grid
(target ${\sim}190$ cells across the larger extent dimension), stamping every
segment cell-by-cell with a one-cell radius so long contours are continuous
in raster space. The remaining cells are solved by Gauss-Seidel relaxation of
the Laplace equation with contour cells held fixed (150--600 iterations,
scaled to grid size) -- the heat-equation smoothing that produces the
smoothest surface consistent with the contour constraints. A BFS
distance-to-data field flags cells far from any contour as interpolation-only
areas.

\textbf{Terrain derivatives.} Slope (percent) is computed by central
differences on the raw surface; local relief (elevation minus the mean of a
$7\times7$ window) identifies natural depressions. OpenCV was considered for
this stage but was not needed: all operations are array kernels over the DEM,
and implementing them directly keeps the pipeline isomorphic (the same code
runs for both modes) and dependency-free.

\subsection{Catchment Area Delineation}
Catchments are derived hydrologically, not buffered geometrically:

\begin{enumerate}
  \item \textbf{Sink filling.} An iterative Planchon--Darboux-style pass
  raises any cell that has no lower neighbour to (lowest neighbour +
  0.001\,m), repeated (max 40 passes) until the surface drains. The number
  of filled cells is reported in the response (\texttt{sinksFilled}).
  \item \textbf{D8 flow direction.} Each cell drains to its steepest
  downslope neighbour among the 8, using the filled surface and true
  neighbour distances ($\sqrt{2}\,c$ for diagonals).
  \item \textbf{Flow accumulation.} Cells are processed high-to-low and each
  adds its accumulated count downstream, so
  \texttt{acc}$_i$ = number of upstream cells draining through cell $i$
  ($O(n \log n)$ for the sort, $O(n)$ for the sweep).
  \item \textbf{Strahler stream order.} Computed on the same high-to-low
  schedule: leaves are order~1, and a cell receiving two or more equal-maximum
  inflows takes order${+}1$ (capped at 8).
  \item \textbf{Upstream catchment.} For a candidate cell,
  \texttt{upstreamCells} walks the reversed D8 graph (BFS over neighbours
  whose flow direction lands in the set), giving the exact contributing cell
  set; \texttt{catchmentArea} = cell count $\times$ cell area. The catchment
  outline is traced from the mask's boundary edges and simplified with a
  ring-aware Douglas--Peucker variant (closed rings are split at the vertex
  farthest from the anchor before simplification -- naive Douglas--Peucker
  degenerates on closed rings and erases them).
\end{enumerate}

\textbf{The village-pond band.} A hard filter restricts candidates to
outlets whose upstream catchment is between 1{,}000 and 6{,}000\,m\textsuperscript{2}:
smaller catchments cannot fill a pond in a typical monsoon; larger ones
describe streams that need check dams, not dug ponds. Within the band, the
scoring term \texttt{bandSizeScore} peaks mid-band (${\sim}3{,}500$\,m\textsuperscript{2}),
so mid-sized catchments outrank band-edge cases. Candidate selection then
applies non-maximum suppression (minimum separation $\approx$
$\min(w,h)/9$ cells) so the returned sites spread across the terrain instead
of clustering on one flow line.

\textbf{Assumptions and limitations.} D8 assigns each cell a single flow
direction, which can parallelise flow across planar slopes and occasionally
picks spurious diagonal routes; the Planchon--Darboux filling adds a nominal
0.001\,m gradient that slightly perturbs accumulation near flats; and the
catchment of a candidate near the grid edge is truncated by the extent.
Catchments are additionally clamped to a physically plausible maximum
(circle area $\times$ 4/$\pi$ bounding-square overshoot).

\subsection{Rainfall Data Integration}
Rainfall is fetched through a three-tier provider chain so analysis always
completes regardless of network availability:
\begin{enumerate}
  \item \textbf{Open-Meteo ERA5 archive} (primary): daily
  \texttt{precipitation\_sum} spanning the three complete water years ending
  last month. Daily values are summed per calendar month and averaged across
  years, giving 12 monthly means and their annual total.
  \item \textbf{Open-Meteo live forecast} (secondary): last 92 days of
  precipitation, annualised by multiplying monthly sums by $365.25/92$.
  Used when the archive endpoint is rate-limited or unreachable.
  \item \textbf{FAO CLIMWAT regional climatology} (offline fallback): a
  hard-coded latitude-band lookup table derived from FAO CLIMWAT long-term
  averages, covering four Indian geographic zones. No network call is
  required; the provider label in the response makes the source explicit.
\end{enumerate}
The UI draws the monthly profile as a bar chart, and every candidate's
explanation quotes the annual figure and its source.

\textbf{Runoff volume estimation.} Annual runoff into a candidate uses the
rational-method volume form
$R = P \cdot A \cdot C$ with $P$ the annual rainfall (m), $A$ the upstream
catchment area (m\textsuperscript{2}), and a fixed runoff coefficient
$C = 0.45$ (a deliberate, documented constant; the input box was removed from
the UI to keep the model honest about its one-parameter simplicity). Each
candidate also reports \texttt{fillRatio} = $R$ / terrain-derived storage, and
a caution is raised when the annual runoff cannot reasonably fill the basin.

\textbf{Pond depth and storage recommendation.} For each candidate,
\texttt{pondForCatchment} grows the pond footprint from the outlet over the
filled surface at the design depth (default 3\,m), traces the footprint
outline, and computes storage as the volume enclosed
(\texttt{storage\_m3}), the mean depth (storage / footprint area), and the
fill ratio against annual runoff.

\textbf{Hard exclusions.} Candidates whose location falls inside a mapped OSM
water polygon are removed unconditionally (\texttt{pointInWater} over closed
rings, including multipolygon relations assembled from member ways).
Proximity flags (within 50\,m of a building, 30\,m of a road) feed the
written reasons and cautions rather than the score.

\textbf{Candidate count.} The number of returned candidates scales with the
study radius:
$\texttt{wantCandidates} = \operatorname{round}(\min(10,\ \max(5,\ 5 + (r - 500)/300)))$
-- 5 for small circles (e.g.\ $r=300$\,m), 7 at $r=1{,}000$\,m, and 10 from
$r=2{,}000$\,m upward. The engine over-generates a $2\times$ pool so that
water exclusions do not thin the list, then slices to the wanted count.

%% =====================================================================
\section{Implementation}
\label{sec:implementation}

\subsection{Backend and API Design}
The backend (\texttt{backend/server.ts}) is a single Express application
with three JSON endpoints and static SPA serving; see
Table~\ref{tab:endpoints}. Design decisions:

\begin{itemize}
  \item \textbf{Plain REST with JSON bodies.} Two analysis resources mirror
  the two input modes; parameters are snake\_case in request bodies and
  camelCase in options, matching the shared TypeScript types.
  \item \textbf{Validation first.} Coordinates, radius bounds
  (100--20{,}000\,m) and upload type/size (only \texttt{.kml}/\texttt{.kmz},
  $\leq$40\,MB, KMZ archives must contain a \texttt{.kml}) are checked before
  any work; failures return HTTP 400 with a specific human-readable message.
  \item \textbf{Errors that tell the truth.} When a mandatory external
  dataset fails, the API returns HTTP 502 with the provider's error and an
  explicit note that no fallback results are generated. When OSM is
  unreachable, the analysis still succeeds but
  \texttt{waterExclusion:\ \{applied:\ false,\ note:\ ...\}} says so.
  \item \textbf{KMZ handling.} Uploads are held in memory; KMZ archives are
  unzipped with \texttt{fflate} and the inner KML is parsed; a ZIP signature
  with a wrong extension is rejected explicitly.
  \item \textbf{Authentication: none.} The service performs read-only
  analysis of public data and holds no user state; auth was scoped out for
  the assignment and is listed as future work (Section~\ref{sec:csd-themes}).
\end{itemize}

\begin{table}[h]
  \caption{REST endpoints}
  \label{tab:endpoints}
  \begin{tabular}{p{1.4cm}p{4.0cm}p{6.6cm}}
    \toprule
    Method & Path & Purpose \\
    \midrule
    GET & \small\texttt{/api/health} & Liveness/version probe polled by the UI every 30\,s. \\
    POST & \small\texttt{/api/analyze} & Circle analysis: \texttt{latitude}, \texttt{longitude}, \texttt{radius\_m} (required); optional \texttt{designDepth\_m}, \texttt{maxSlopePct}, \texttt{minCatchmentArea\_m2}, \texttt{maxCandidates}, \texttt{runoffCoefficient}. Returns DEM metadata, rainfall stats, water layers and ranked candidates. \\
    POST & \small\texttt{/api/analyzeContour} & KML/KMZ contour-sheet analysis (multipart field \texttt{file}); optional circle fields restrict candidates to the radius; same response shape plus contour/radius-filter metadata. \\
    GET & \small\texttt{/*} (non-\texttt{/api}) & Serves the built SPA from \texttt{dist/} with history fallback (production). \\
    \bottomrule
  \end{tabular}
\end{table}

\subsection{Frontend and Visualization}
The frontend is a React~19 single-page application built with Vite. The
workspace pairs a map panel with an input/results panel:

\begin{itemize}
  \item \textbf{Map (Leaflet~1.9, \texttt{PondMap.tsx}).} Three base maps --
  Street (OSM), Terrain (OpenTopoMap) and Satellite (Esri World Imagery);
  dashed study-radius circle with centre marker; blue OSM water bodies and
  waterways; uploaded contour polylines; translucent catchment polygons for
  each candidate; and \texttt{CS-XX} chips on candidate sites. Clicking a
  chip or catchment opens a popup with the candidate's score, status,
  confidence, factors and explanation. Layers are drawn only after the map
  is ready (eliminating an initialization race), and the Leaflet stylesheet
  is imported explicitly (its absence was the cause of mis-rendered tiles).
  \item \textbf{Workspace panel (\texttt{Workspace.tsx}).} Study-input card
  (latitude, longitude, radius), contour-analysis card (KML/KMZ upload with
  optional circle), export card (JSON download of the full response), the
  rainfall bar chart, and the selected candidate's scoring-factor breakdown
  with weights, values and written detail lines.
  \item \textbf{Health pill.} A status indicator polls
  \texttt{/api/health} every 30\,s so the planner always knows whether the
  analysis backend is reachable.
  \item \textbf{Report view (\texttt{Report.tsx}).} A printable summary of
  the current analysis (DEM, rainfall, water exclusion status, candidate
  table).
\end{itemize}

Figure~\ref{fig:ui} shows the workspace: the candidate map with contour base
map and \texttt{CS-XX} chips (left) and the input/analysis/export cards
(right). Figure~\ref{fig:results} in Section~\ref{sec:results} shows the
street base map with water bodies and the legend.

\begin{figure}[h]
  \centering
  \IfFileExists{assets/workspace.png}%
    {\includegraphics[width=\linewidth]{assets/workspace.png}}%
    {\fbox{\parbox[c][6cm][c]{0.92\linewidth}{\centering
      \textit{Screenshot placeholder: save the workspace screenshot as}\\
      \texttt{report/assets/workspace.png}\\ \textit{and recompile.}}}}
  \caption{Application interface: map overlay (catchment polygons, water
    bodies, CS-XX candidate chips) beside the study-input, contour-analysis
    and export cards, with the rainfall profile and scoring-factor
    breakdown.}
  \label{fig:ui}
  \Description{Screenshot of the PondSite workspace. Left: a Leaflet map
    with candidate chips CS-01 through CS-NN, translucent catchment
    polygons, and blue OSM water bodies. Right: the STUDY INPUT card with
    latitude, longitude and radius fields, the CONTOUR ANALYSIS card with a
    KML/KMZ upload control, the EXPORT card, a monthly rainfall bar chart,
    and a candidate scoring-factor breakdown.}
\end{figure}

\subsection{Database and Storage}
There is deliberately no server-side database. Spatial data, rainfall
records and computed recommendations live only for the duration of a request
and are returned to the client in full. Persistence responsibilities are
distributed as follows: (i)~the \emph{browser} keeps the most recent
analysis in \texttt{localStorage} so a page reload restores the workspace;
(ii)~the \emph{export} card writes the complete JSON response (DEM metadata,
rainfall series, water layers, every candidate with factors and
explanations) to a file for records; (iii)~uploaded KML/KMZ files are held
in memory by \texttt{multer}, parsed, and discarded -- nothing from the
upload path is written to disk. This keeps the service stateless (any
instance can serve any request), avoids stale-data bugs, and matches the
no-fallback data philosophy: stored results would eventually contradict what
the live APIs return.

%% =====================================================================
\section{CSD Themes and Topics Applied in the Project}
\label{sec:csd-themes}
Table~\ref{tab:csd-themes} maps the CSD course themes to concrete components
of PondSite, with the design rationale for each. Themes that were not
implemented (load balancing, database indexing) are omitted rather than
claimed.

\begin{table*}[t]
  \caption{Mapping of core CSD themes/topics to their use in this project}
  \label{tab:csd-themes}
  \begin{tabular}{p{2.2cm}p{4.0cm}p{5.6cm}}
    \toprule
    CSD Theme/Topic & Where used in the project & Justification / design rationale \\
    \midrule
    API Design (REST) & \texttt{POST /api/analyze}, \texttt{POST /api/analyzeContour}, \texttt{GET /api/health} (\texttt{server.ts}) & REST with JSON bodies keeps the contract small and type-shared with the frontend; strict validation returns 400/502 with actionable messages instead of silent failures. \\
    Concurrency / Asynchronous Processing & \texttt{Promise.allSettled} over DEM + rainfall + OSM fetches (\texttt{analyze-circle.ts}); \texttt{Promise.any} racing 7 Overpass mirrors (\texttt{overpass.ts}) & Sequential fetching made analyses take minutes when one API was slow; parallel fan-out bounds end-to-end latency by the slowest mandatory source, and mirror racing removes single-host stalls. \\
    Error Handling and Resilience & 3-retry backoff per DEM batch (\texttt{elevation.ts}); explicit \texttt{waterExclusion:\ \{applied:\ false\}}; no-fallback policy for terrain/rainfall & The system degrades gracefully: optional data is reported as missing, mandatory failures are honest 502s -- never invented climatology or synthetic terrain. \\
    Algorithms and Complexity & Sink fill $O(40n)$, D8 + accumulation $O(n \log n)$, upstream BFS $O(n)$, Strahler ordering, Douglas--Peucker ring simplification (\texttt{contour.ts}) & D8 was chosen over multi-flow-direction for exact, reproducible contributing-area sets at $n \le 4{,}096$ cells; the sort-based accumulation schedule guarantees each cell is settled before its downstream. \\
    Design Patterns & Provider-fallback strategy (\texttt{elevation.ts}); shared pipeline modules reused by both analysis modes & The elevation module encapsulates provider order and retries behind one \texttt{fetchElevations} call; both circle and contour modes reuse the identical hydrology core. \\
    Microservices vs.\ Monolith & Deliberate modular monolith (\texttt{backend/lib/*}) & A second service would add deployment and latency overhead with no scaling benefit for a single-tenant analyser; module boundaries (\texttt{contour}, \texttt{elevation}, \texttt{rainfall}, \texttt{overpass}) keep extraction cheap if requirements grow. \\
    Caching & Browser/Leaflet tile caching; analysis results kept client-side in \texttt{localStorage} & Server-side response caching was considered and rejected: rainfall windows roll monthly and OSM changes continuously, so cached analyses would silently diverge from live data. \\
    Testing Strategy & \texttt{tsc --noEmit} strict typecheck, production build, and live end-to-end API verification against real coordinates and the sample \texttt{contours\_1m.kml} & Type-level checks plus real-data runs verify the no-fallback contract (all values traceable to providers) better than mocked unit tests of the data layer. \\
    Version Control / CI-CD & Git-tracked repository with README-documented run/build scripts (\texttt{npm run dev}, \texttt{npm run build}, \texttt{npm start}) & Single-developer course project: conventional commits and a reproducible two-command run kept the tree reviewable; a CI pipeline is future work. \\
    Authentication / Authorization & Not implemented & The service exposes read-only analysis of public data and stores no user state; adding auth was judged out of scope and is listed as future work. \\
    \bottomrule
  \end{tabular}
\end{table*}

Three themes shaped the system most. \emph{Concurrency}: the original
sequential pipeline fetched terrain, then rainfall, then OSM; a single slow
provider stalled the whole analysis for minutes. Restructuring to
\texttt{Promise.allSettled} (with mandatory/optional distinction on the
results) and racing all Overpass mirrors with \texttt{Promise.any} brought
typical end-to-end time under a minute. \emph{Error handling and
resilience}: the assignment's hardest engineering decision was refusing to
fall back. Early prototypes invented plausible numbers when APIs failed,
which silently invalidates every downstream recommendation; the shipped
policy is that terrain and rainfall failures abort with honest errors, and
OSM unavailability is reported structurally in the response.
\emph{Algorithms and complexity}: D8 over a capped $64\times64$ grid keeps
the entire hydrology below a few million operations per analysis, which is
what makes live, interactive catchment delineation possible at all.

%% =====================================================================
\section{Results and Evaluation}
\label{sec:results}
The system was verified end-to-end against live providers (all numbers below
are from actual runs, not fixtures).

\textbf{Circle analysis (500\,m radius at 21.2591\,N, 81.3002\,E, Chhattisgarh, India).}
The analysis returned 5 candidates -- the correct value for the radius band
formula at $r = 500$\,m -- every one with an upstream catchment inside the
1{,}000--6{,}000\,m\textsuperscript{2} band. The DEM source reported was
SRTM 30\,m via OpenTopoData (Open-Meteo was rate-limited with HTTP 429 at
the time -- the fallback chain behaved as designed). ERA5 rainfall
summed to 1{,}123\,mm/year with the expected monsoon-concentrated monthly
profile. Overpass was reachable, water exclusion applied, and existing water
bodies and waterways were drawn on the map.

\textbf{Contour analysis (uploaded 1\,m contour KML, extent
21.2398--21.2636\,N, 81.2814--81.3126\,E).} The file produced 14 candidates,
all with catchments of 3{,}201--3{,}784\,m\textsuperscript{2} -- inside the
band and, notably, clustered near the band mid-point that
\texttt{bandSizeScore} favours, indicating the DEM interpolation and hydrology
are behaving consistently. Re-running the same file with a 1{,}200\,m study
circle kept 4 of the 14 candidates, confirming the radius filter.

\textbf{Explainability.} Every candidate carries its four weighted factors
(flow convergence 0.36, slope 0.28, natural depression 0.21, data support
0.15) with the underlying numbers in the detail text, plus written reasons,
cautions, and a plain-language explanation quoting catchment area, annual
runoff at 1{,}123\,mm and $C = 0.45$, and terrain-derived storage at 3\,m
depth. Status badges (RECOMMENDED / CONDITIONAL / REJECTED) and confidence
(HIGH/MEDIUM/LOW) are derived from the same published weights, so a reviewer
can recompute any score by hand.

Figure~\ref{fig:results} shows the results overlay on the street base map
with water bodies and the legend.

\begin{figure}[h]
  \centering
  \IfFileExists{assets/candidates-map.png}%
    {\includegraphics[width=\linewidth]{assets/candidates-map.png}}%
    {\fbox{\parbox[c][6cm][c]{0.92\linewidth}{\centering
      \textit{Screenshot placeholder: save the candidates-map screenshot as}\\
      \texttt{report/assets/candidates-map.png}\\ \textit{and recompile.}}}}
  \par\medskip
  \IfFileExists{assets/street-map.png}%
    {\includegraphics[width=\linewidth]{assets/street-map.png}}%
    {\fbox{\parbox[c][6cm][c]{0.92\linewidth}{\centering
      \textit{Screenshot placeholder: save the street-map screenshot as}\\
      \texttt{report/assets/street-map.png}\\ \textit{and recompile.}}}}
  \caption{Example results overlay. Top: candidates over the contour base
    map with catchment polygons and \texttt{CS-XX} chips. Bottom: street
    base map with OSM water bodies, waterways, the legend and candidate
    chips.}
  \label{fig:results}
  \Description{Two screenshots of the PondSite map. The first shows
    candidates over a contour base map with translucent catchment polygons
    and CS-numbered chips. The second shows the same area on the street base
    map with blue water bodies and waterways, a legend, and CS-01 through
    CS-20 candidate chips.}
\end{figure}

\subsection{Performance}
\textbf{End-to-end latency.} A full circle analysis (DEM sampling + rainfall
+ OSM + hydrology + scoring) typically completes in well under a minute on a
residential connection; the wall-clock is dominated by DEM batch fetches
(up to 41 batches of 100 points), which the retry/backoff policy paces to
respect provider rate limits. Before the parallelisation described in
Section~\ref{sec:csd-themes}, the same pipeline took minutes whenever one
provider was slow; the fan-out now bounds latency by the slowest
\emph{mandatory} source.

\textbf{Component behaviour.} Overpass responses are bounded by the 60\,s
timeout across seven raced mirrors; DEM and rainfall requests carry 30\,s
timeouts with up to three retries each. Local operations are fast relative
to the network: the capped hydrology grid ($\le 4{,}096$ cells) and the
Laplace relaxation for contour grids ($\le$600 iterations over ${\sim}190$
columns) run in tens of milliseconds on a laptop-class CPU. The health
endpoint responds sub-second locally and is polled every 30\,s. Load testing
under concurrent users was not performed (single-tenant course deployment);
the stateless design means horizontal scaling is trivial if ever needed.

%% =====================================================================
\section{Discussion and Limitations}
\label{sec:discussion}
\textbf{What worked well.} The no-fallback policy proved its worth during
development: when Open-Meteo rate-limited the machine (HTTP 429) and several
Overpass mirrors were network-blocked, the system either used its alternate
provider (SRTM via OpenTopoData), reported water exclusion as not-applied, or
failed loudly -- it never produced plausible-looking fiction. The
catchment-band hard filter (1{,}000--6{,}000\,m\textsuperscript{2}) also did
exactly what it was designed to do: across both verification runs, every
single returned candidate fell inside the band.

\textbf{Data-quality issues.} The DEM providers serve 90\,m (Copernicus
GLO-90) or 30\,m (SRTM) native resolution; sampling at ${\sim}20$\,m cannot
recover detail finer than the source, so micro-depressions and narrow
drainage lines are smoothed away. In contour mode, Laplace interpolation
between 1\,m contours is smooth but optimistic: it cannot invent valley
thalwegs between widely spaced lines, and the result depends on contour
coverage density. ERA5 reanalysis rainfall is gridded model data, not a
gauge record -- useful for relative comparisons, less exact for absolute
volumes.

\textbf{Model simplifications that affect accuracy.} The runoff coefficient
is a fixed $C = 0.45$: real coefficients vary with land cover, soil and
antecedent moisture, so absolute runoff volumes should be read as screening
estimates. Point-in-water exclusion uses OSM geometry as-is (no safety
buffer around seasonal tanks), and OSM completeness varies by region --
candidates are never excluded by water OSM does not know about. D8
single-direction routing can mis-route flow across flat or planar terrain
despite sink filling. The pond storage model assumes vertical-sided
excavation from the filled surface at the design depth and ignores seepage
and evaporation in the fill-ratio check.

\textbf{Engineering limitations.} Large KMZ uploads are memory-parsed with a
40\,MB cap; heavier contour sheets would need streaming. There is no
server-side caching, so repeated analyses of the same circle re-pay the
network cost -- correct for freshness, expensive for iteration. Finally,
Open-Meteo's rate limiting was observed to be IP-based; institutional
deployments behind shared egress IPs will hit the SRTM fallback more often,
which works but is slower and coarser.

%% =====================================================================
\section{AI Tool Usage Declaration}
\label{sec:ai-usage}
AI tools were used during this project in accordance with the assignment's
LLM usage policy. Specifically, the Codebuff AI coding agent (running the
GLM model) was used for: (i)~scaffolding and refactoring project files
(backend modules, React components, configuration); (ii)~debugging -- for
example diagnosing a missing Leaflet CSS import that corrupted map tiles, a
non-monotonic initialisation race in map layer drawing, and a
Douglas--Peucker degeneracy on closed catchment rings; (iii)~formatting and
restructuring this report from the provided template; and (iv)~drafting
README documentation.

All AI-assisted code was reviewed, understood, tested, and where necessary
modified by the author before submission: the full project typechecks under
\texttt{tsc --noEmit} (strict), builds cleanly with Vite, and every
reported result in Section~\ref{sec:results} was reproduced against live
APIs by the author. Algorithm descriptions in Section~\ref{sec:methodology}
were written by reading and verifying the final code, not generated
blindly. No AI-generated content is presented as an unreviewed artifact.


%% =====================================================================
\appendix

\section{Source Code and Repository}
%% TODO: replace with your actual repository URL before submission.
The complete source code is available at
\texttt{<your-repository-url>} (submission archive \texttt{pondsite.zip}).
The folder structure:

\begin{verbatim}
pond-planner-pro/
  backend/
    server.ts                 Express API + static SPA serving
    lib/
      analyze-circle.ts       circle pipeline (DEM -> hydrology -> scoring)
      analyze-contour.ts      contour engine (parsing, DEM, ranking)
      analyze-contour-wrapper.ts  KML/KMZ path (full-extent mode, no radius)
      contour.ts              DEM/hydrology primitives (Laplace fill, D8, catchments)
      elevation.ts            4-tier DEM: Open-Meteo -> OpenTopoData -> Open-Elevation -> synthetic
      rainfall.ts             3-tier fallback: ERA5 -> forecast -> FAO climatology
      history.ts              JSON file store (last 50 analyses, auto-saved)
      overpass.ts             OSM water/buildings/roads + exclusion helpers
      geo.ts, types.ts        shared geometry + API types
  src/
    App.tsx                   decoupled circle + KML state, health polling, nav
    components/PondMap.tsx    Leaflet map (3 base maps, catchment labels, popups)
    pages/Workspace.tsx       two independent tabs: Circle Analysis / KML Analysis
    pages/Report.tsx          two independent report sections
    pages/History.tsx         history browser with map replay
    index.css                 Leaflet CSS + catchment label styles
  deploy/
    lb/main.go                Go load balancer (round-robin, health-check, SPA serve)
    ecosystem.config.cjs      PM2: 1 Go LB (port 3205) + 4 Node backends (3001-3004)
    deploy.sh                 one-shot server deploy script
  report/                     this report + screenshot assets
  contours_1m.kml             sample 1 m contour file for testing
\end{verbatim}

\textbf{Live deployment.} The system is hosted at
\url{http://10.1.75.51:3205}. The production stack runs four
Node.js backend instances (ports 3001--3004) managed by PM2 for
automatic restart, with a custom Go load balancer on port~3205 distributing
traffic round-robin across them. The sample file
\texttt{contours\_1m.kml} (included in the archive) exercises the KML
upload path directly in the browser.

\end{document}
\endinput
